\documentclass[14pt]{extarticle}
\usepackage{preamble}
\geometry{letterpaper, landscape, margin=0.5in}
\renewcommand{\answer}[1]{}
\setlength{\parskip}{4pt}

\begin{document}
\pagestyle{empty}

\daybanner{Unit 2 \textperiodcentered\ Topic 2.6 \textperiodcentered\ THE PROBLEM}{Which Way Does the Condition Point?}

\noindent\begin{minipage}[t]{0.57\linewidth}
\vspace{0pt}
Suppose a school screens 1{,}000 students for a hypothetical condition. Exactly $5\%$ have the condition. The test has $90\%$ sensitivity: among students with the condition, $90\%$ test positive. Its false-positive rate is $10\%$: among students without the condition, $10\%$ test positive. The implied counts are shown.\par\smallskip \textbf{Mina:} ``The test catches $90\%$ of students who have the condition, so a positive result means a student has a $90\%$ chance of having it.''\par \textbf{Luis:} ``We need to count everyone who tested positive before finding the chance of having the condition.''\par
\end{minipage}\hfill
\begin{minipage}[t]{0.40\linewidth}
\vspace{0pt}
\begin{center}
\textbf{Screening results (counts)}\par\smallskip
\begin{tabular}{|l|c|c|c|}
\hline
\textbf{Status} & \textbf{Pos.} & \textbf{Neg.} & \textbf{Total} \\ \hline
\textbf{Condition} & 45 & 5 & 50 \\ \hline
\textbf{No condition} & 95 & 855 & 950 \\ \hline
\textbf{Total} & 140 & 860 & 1{,}000 \\ \hline
\end{tabular}
\end{center}
\end{minipage}

\begin{firsttakebox}{FIRST TAKE --- one sentence before you work the parts}
Does a positive result mean a $90\%$ chance of having the condition? Choose a claim and cite one table count.
\end{firsttakebox}

\begin{description}[leftmargin=1.15in, labelwidth=1in, itemsep=2pt, topsep=2pt]
  \item[\dokbadge{1}~(a)] For $P(\text{condition}\mid\text{positive})$, identify the numerator and denominator counts.
  \item[\dokbadge{2}~(b)] Find $P(\text{positive}\mid\text{condition})$ and $P(\text{condition}\mid\text{positive})$. Show each count fraction and percent.
  \item[\dokbadge{3}~(c)] Is Mina's claim supported? Use the positive-test counts for students with and without the condition. Explain how reversing given changes the group in the denominator. Write 3--4 sentences.
\end{description}

\vfill
\textbf{Self-paced bonus problem. Work (a)--(c) from this sheet; turn it in whenever you finish.}

\end{document}
